\section{Conclusion}
\label{sec:conclusion}

\subsection{Summary}

The \flip{} proposal is the simplest valid redistricting perturbation: move
one boundary tract from its current district to an adjacent district, accept
if the result is contiguous and population-balanced.
Its properties follow directly from this simplicity:

\paragraph{Fast.}
Each step runs in $O(\deg)$ amortised time with incremental boundary set
maintenance.
Ten thousand steps take under 3 seconds for any state in our implementation.

\paragraph{High acceptance.}
Boundary tracts are near-balanced by selection; contiguity failure is rare.
Acceptance rates of $40$--$70\%$ are typical, giving $4{,}000$--$7{,}000$
distinct visited plans per $10{,}000$-step run.

\paragraph{Locally confined.}
Each \flip{} step changes exactly one tract.
The chain cannot tunnel through topological obstacles (Theorem~\ref{thm:reachability}),
so it characterises only the local neighbourhood of the starting plan.
This is a feature for sensitivity analysis and a bug for ensemble sampling.

\paragraph{Slow mixing.}
Moving from any starting plan to a fundamentally different plan requires
$O(\text{district size})$ accepted steps at minimum (Theorem~\ref{thm:slow-mix}).
In practice, the chain does not mix in $10{,}000$ steps for any real-world
state; it traces a random walk in the local neighbourhood.

\subsection{Usage Guidelines}

The \flip{} chain should be used as follows in the redistricting workflow:

\begin{enumerate}
  \item \textbf{After plan generation.} Generate the plan using \cs{} (or
        another global search method). Do not use \flip{} to generate the
        plan; use it only to certify the generated plan.
  \item \textbf{Run the $p$-sweep.} Execute five \flip{} runs at
        $p \in \{0.0, 0.25, 0.5, 0.75, 1.0\}$.
        Record $\rho_{\mathrm{seat}}$ for each.
  \item \textbf{Disclose the results.} Report \texttt{visited\_count},
        \texttt{selected\_plan\_rank}, $\sigma_{\ec}$, and
        $\rho_{\mathrm{seat}}$ in the legislative findings.
        If any $p$-value produces a different seat outcome, explain the
        choice of percentile.
  \item \textbf{Store the manifest.} Record \texttt{base\_seed},
        \texttt{chain\_seed\_0}, \texttt{visited\_count}, and
        \texttt{selected\_plan\_rank} for verifier reproduction.
\end{enumerate}

\subsection{Open Questions}

\begin{enumerate}
  \item \textbf{Optimal step count.} Is $F = 10{,}000$ steps sufficient to
        characterise the local neighbourhood for all states?
        For large states ($n > 5{,}000$) or states with high district counts
        ($k > 20$), a larger $F$ may be warranted.
        A stopping criterion based on the saturation of the visited-plan count
        would provide a principled answer.

  \item \textbf{Weighting the acceptance.} The current \flip{} implementation
        uses a hard-constraint acceptance criterion (accept if contiguous and
        balanced, reject otherwise).
        A soft Metropolis criterion that penalises large EC changes could
        concentrate the chain near the starting plan's EC level, making the
        $p$-sweep more informative about plan quality.

  \item \textbf{Multi-tract flip.} Allowing simultaneous moves of $r > 1$
        adjacent boundary tracts would bridge topological obstacles and
        expand the reachable plan set, at the cost of $O(r!)$ more proposals
        per step.
        This generalisation would interpolate between \flip{} ($r = 1$)
        and \recom{} ($r \approx n/k$).

  \item \textbf{Formal partisan stability bounds.} Given a census-tract graph
        with known geometry (e.g., compact convex districts), can we derive
        a lower bound on $\rho_{\mathrm{seat}}$ under \flip{}?
        Such a bound would let courts distinguish geometric constraint
        (high stability for topological reasons) from deliberate placement
        (artificially high stability due to gerrymandering).
\end{enumerate}

\subsection{Relation to the Full Toolkit}

The redistricting toolkit is now complete at three perturbation levels.
Table~\ref{tab:comparison} (Section~\ref{sec:comparison}) summarises when
to use each method.
The standard workflow is:
\begin{center}
  \cs{} (generate) $\to$ \flip{} (certify) $\to$ \be{} (structure audit)
  $\to$ ensemble (global audit).
\end{center}
Each step adds a different type of evidence.
\flip{} occupies an essential niche: it is the only method that directly
characterises the plan's local neighbourhood at the tract level, and it is
the only method fast enough to run interactively (under 10 seconds) during
the plan-drafting process.
